\documentclass[ASNA,twocolumn]{USG} %twocolumn
\usepackage{anyfontsize} %

\usepackage{colortbl}   % For coloring table lines
\usepackage{xcolor}     % For color definitions

\usepackage{steinmetz}
\graphicspath{{./images/}}
\specialissue{}

\articletype{ORIGINAL ARTICLE}%
\subarticletype{}

\received{}
\revised{}
\accepted{}
\journal{AECAI-PRiSM 2026}
\volume{}
\copyyear{}
\startpage{1}
\articledoi{}

%\raggedbottom

%%
%\renewcommand\ShowFrameLinethickness{0.15pt}
%\renewcommand*\ShowFrameColor{\color{red}}


\begin{document}
\title{Language-based Augmentation of a Virtual Clinical Avatar with LoRA parameter Fine-Tuning}
%\subtitle{Subtitle}
\transtitle{Language-based Augmentation of a Virtual Clinical Avatar with LoRA parameter Fine-Tuning}
\subtranstitle{}
\author[1]{Anonymous Author 1}

\author[1]{Anonymous Author 2}

\authormark{Anonymous \textsc{et al.}}
\titlemark{Language-based Augmentation of a Virtual Clinical Avatar with LoRA parameter Fine-Tuning}

\address[1]{\orgdiv{Department withheld for review, }\orgname{Institution withheld for review, }%
\orgaddress{\state{State/Province withheld, }\country{Country withheld}}}

\corres{Corresponding author withheld for double-blind review.}

\editor{\textbf{Academic Editor:} \textbf{Guest Editor:} }

\presentaddress{}

\fundingInfo{No funding to disclose.}

\keywords{parameter-efficient fine-tuning, low-rank adaptation, pediatric clinical NLP, virtual reality, on-device inference, small language models}

\transkeywords{parameter-efficient fine-tuning, low-rank adaptation, pediatric clinical NLP, virtual reality, on-device inference, small language models}

\abstract[ABSTRACT]{Pediatric patients navigating hospital environments require age-appropriate, institutionally grounded information that general-purpose AI systems are not designed to provide. We present an age-targeted fine-tuning framework for a virtual hospital guide deployed within a digital recreation of **** Hospital (Location withheld for anonymity), specialized for children aged 5--11 and adolescents aged 12--18. Using Low-Rank Adaptation (LoRA) on 4-bit quantized Llama 3.2 Instruct models (1B and 3B) running on an Apple Mac Mini M4 (16~GB), a deliberate constraint modeling resource-limited clinical deployment, we train two adapter configurations: Standard ($r=8$, 16 layers) and Fast ($r=4$, 8 layers), across both model sizes and age groups for 8 training runs total. Training data consists of approximately 560 synthetic examples per adapter (1123 total across six topic categories) generated by Claude Opus 4.6 using **** Hospital's public patient materials as grounding artifacts, following a LIMA-inspired curation approach \cite{zhou2023lima}. Fine-tuned adapters substantially outperform zero-shot base models on readability (72--84\% of 5--11 responses meeting a Flesch-Kincaid $\leq 7.0$ grade-level target vs.\ 12--14\% for base models) and inter-role style separation (TF-IDF + logistic regression classifier accuracy 0.89--0.94 vs.\ 0.66--0.70). A configuration-ordering crossover is observed: the Standard adapter yields better results on the 3B model while the Fast adapter performs better on the 1B model, suggesting optimal adapter capacity is model-size-dependent. A subsequent controlled rank sweep (fixed \texttt{num\_layers~=~8}, $r \in \{2, 4, 8, 16\}$, two seeds, four architectures) identifies rank as the operative factor via capacity regularization, finding that transformer depth, not total parameter count, predicts the optimal rank regime. Cross-architecture evaluation on Qwen~2 0.5B and SmolLM2 360M confirms the crossover is not Llama-specific; SmolLM2 Standard achieves 84\% FK pass rate and 0.950 classifier accuracy at 0.81\,s average latency despite only 360\,M parameters. A subsequent Qwen 2.5 family sweep (0.5B, 1.5B, 3B) finds Fast outperforming Standard uniformly across all three sizes, with no crossover within the family, consistent with the small-model rank regime prediction, with Qwen 2.5 1.5B Fast setting the highest inter-role classifier accuracy of any adapter evaluated (0.980). A final layer sweep (fixed rank=4, \texttt{num\_layers} $\in \{4, 8, 16\}$, Llama 1B/3B, two seeds) confirms that layer depth is an independent secondary contributor to the crossover: at fixed rank, layers=16 outperforms layers=8 by 4--7\% FK pass rate, while small models peak at layers=4. These results demonstrate that compact, parameter-efficient adaptation on consumer hardware is feasible for pediatric clinical NLP, and that both adapter rank and layer depth jointly determine optimal configuration, with architecture depth predicting the rank regime.
}


 \transabstract[transABSTRACT]{Pediatric patients navigating hospital environments require age-appropriate, institutionally grounded information that general-purpose AI systems are not designed to provide. We present an age-targeted fine-tuning framework for a virtual hospital guide deployed within a digital recreation of **** Hospital (Location withheld for anonymity), specialized for children aged 5--11 and adolescents aged 12--18. Using Low-Rank Adaptation (LoRA) on 4-bit quantized Llama 3.2 Instruct models (1B and 3B) running on an Apple Mac Mini M4 (16~GB), a deliberate constraint modeling resource-limited clinical deployment, we train two adapter configurations: Standard ($r=8$, 16 layers) and Fast ($r=4$, 8 layers), across both model sizes and age groups for 8 training runs total. Training data consists of approximately 560 synthetic examples per adapter (1123 total across six topic categories) generated by Claude Opus 4.6 using **** Hospital's public patient materials as grounding artifacts, following a LIMA-inspired curation approach \cite{zhou2023lima}. Fine-tuned adapters substantially outperform zero-shot base models on readability (72--84\% of 5--11 responses meeting a Flesch-Kincaid $\leq 7.0$ grade-level target vs.\ 12--14\% for base models) and inter-role style separation (TF-IDF + logistic regression classifier accuracy 0.89--0.94 vs.\ 0.66--0.70). A configuration-ordering crossover is observed: the Standard adapter yields better results on the 3B model while the Fast adapter performs better on the 1B model, suggesting optimal adapter capacity is model-size-dependent. A subsequent controlled rank sweep (fixed \texttt{num\_layers~=~8}, $r \in \{2, 4, 8, 16\}$, two seeds, four architectures) identifies rank as the operative factor via capacity regularization, finding that transformer depth, not total parameter count, predicts the optimal rank regime. Cross-architecture evaluation on Qwen~2 0.5B and SmolLM2 360M confirms the crossover is not Llama-specific; SmolLM2 Standard achieves 84\% FK pass rate and 0.950 classifier accuracy at 0.81\,s average latency despite only 360\,M parameters. A subsequent Qwen 2.5 family sweep (0.5B, 1.5B, 3B) finds Fast outperforming Standard uniformly across all three sizes, with no crossover within the family, consistent with the small-model rank regime prediction, with Qwen 2.5 1.5B Fast setting the highest inter-role classifier accuracy of any adapter evaluated (0.980). A final layer sweep (fixed rank=4, \texttt{num\_layers} $\in \{4, 8, 16\}$, Llama 1B/3B, two seeds) confirms that layer depth is an independent secondary contributor to the crossover: at fixed rank, layers=16 outperforms layers=8 by 4--7\% FK pass rate, while small models peak at layers=4. These results demonstrate that compact, parameter-efficient adaptation on consumer hardware is feasible for pediatric clinical NLP, and that both adapter rank and layer depth jointly determine optimal configuration, with architecture depth predicting the rank regime.
}


\abbr{abbreviations here}

\contributed{}

\dedicated{}

\copyright{This is an open access article under the terms of the \href{Creative Commons Attribution-NonCommercial}{Creative Commons Attribution-NonCommercial} License, which permits use, distribution and reproduction in any medium, provided the
original~work~is~properly cited and is not used for commercial purposes.
\\[5pt]
  ©  2026 The Author(s).}

%\openaccessstatement

\maketitle
%\afterpage{\aftergroup\restoregeometry}

%\onecolumn
\section{Introduction}\label{sec1}

{L}{arge} language models (LLMs) have created new opportunities for domain-specific conversational systems, but deploying them in resource-constrained clinical environments introduces practical constraints that are rarely addressed together: limited memory budgets favor smaller models, domain adaptation requires fine-tuning on specialized data, and real-time interaction imposes strict latency bounds.

This paper presents \textit{VR Hospital Avatar}, an age-targeted conversational guide embedded in a virtual recreation of **** Hospital (Location withheld for anonymity). The system provides age-appropriate hospital information to two patient groups: children aged 5--11 and adolescents aged 12--18, via natural language interaction within a VR environment. All training and inference runs on an Apple Mac Mini M4 (16~GB), a deliberate constraint representing deployment on consumer hardware without cloud infrastructure.

To specialize the system for each age group within this memory budget, we fine-tune Llama 3.2 Instruct models \cite{dubey2024llama3} at 1B and 3B scales using Low-Rank Adaptation (LoRA) \cite{hu2022lora}. Training data is generated synthetically following the LIMA hypothesis \cite{zhou2023lima}: approximately 560 training examples per adapter (1123 total across six topic categories) grounded in **** Hospital's public patient and family materials, rather than a large noisy corpus. We compare two adapter configurations (Standard: $r = 8$, 16 adapted layers; Fast: $r = 4$, 8 adapted layers) across both model sizes and age groups for eight total training runs.

Our evaluation surfaces a \textbf{configuration-ordering crossover}: the Standard adapter yields better readability and style-separation results on the 3B model, while the Fast adapter outperforms on the 1B model. This crossover is consistent across all five readability metrics and the inter-role classification task, and suggests that optimal adapter capacity is model-size-dependent, a dependency not captured by common LoRA guidelines that recommend a fixed rank range without model-size conditioning \cite{biderman2024lora}. Only the Fast-1B configuration meets the 1.0-second real-time latency target for VR interaction, specifically for the 5--11 age group (0.93s average); the 12--18 adapter averages 1.20s.

The contributions of this paper are:
\begin{itemize}
    \item A complete age-adaptive fine-tuning pipeline for pediatric hospital communication, trained and evaluated on consumer hardware.
    \item Empirical evidence of a LoRA configuration-ordering crossover between Llama 3.2 1B and 3B, consistent across readability and style-separation metrics.
    \item Mechanism identification via a controlled rank sweep (32 runs: 4 ranks $\times$ 4 architectures $\times$ 2 seeds), confirming rank-driven capacity regularization and establishing that transformer depth predicts optimal rank regime.
    \item Cross-architecture validation on Qwen~2 0.5B, SmolLM2 360M, and the Qwen 2.5 family (0.5B--3B), confirming the crossover is not Llama-specific and that Qwen 2.5 Fast adapters dominate Standard uniformly across the full capacity ladder, consistent with the small-model rank regime.
    \item Layer sweep evidence (fixed rank=4, \texttt{num\_layers} $\in \{4, 8, 16\}$, Llama 1B/3B, two seeds) that depth is an independent secondary contributor: layers=16 outperforms layers=8 by 4--7\% FK at fixed rank, while small models peak at layers=4, confirming the crossover reflects the joint effect of both rank and depth.
\end{itemize}

\section{Related Work}\label{sec2}

\subsection{Parameter-Efficient Fine-Tuning}

Low-Rank Adaptation (LoRA) \cite{hu2022lora} reduces fine-tuning cost by constraining weight updates to a low-rank decomposition $\Delta W = BA$, leaving the pre-trained weights frozen. QLoRA \cite{dettmers2023qlora} extended this to 4-bit quantized base models via NormalFloat quantization, enabling fine-tuning of 7B--65B models on consumer hardware. Both are now standard in resource-constrained deployment pipelines.

Rank selection in LoRA is typically treated as a task-dependent hyperparameter. Common practitioner guidance, including HuggingFace PEFT defaults and QLoRA \cite{dettmers2023qlora}, recommends rank 8--16 without conditioning on model scale. Biderman et al.~\cite{biderman2024lora} provide the most thorough empirical study to date, recommending rank 16 as a practical default, but this guidance is developed and validated at the 7B--70B scale; rank ordering behavior below 7B is not characterized by existing work. Our work provides empirical evidence that adapter performance ordering can reverse between 1B and 3B models on the same task, a regime not covered by prior rank selection studies.

\subsection{LLMs in Clinical and Health Communication}

LLMs have been applied to patient education, clinical FAQ systems, and virtual health assistants. Pediatric deployment introduces challenges distinct from adult-facing systems: responses must match developmental vocabulary, maintain an emotionally supportive (rather than clinical) register, and enforce hard safety constraints such as never implying a diagnosis. These properties are not delivered by generic instruction-following fine-tuning on a general-purpose model.

A recurring constraint in clinical NLP is the unavailability of naturalistic patient dialogue due to privacy regulations, motivating synthetic data generation. The LIMA hypothesis \cite{zhou2023lima}, which holds that a small set of carefully curated, high-quality examples suffices for behavioral alignment, supports viable training corpora grounded in institutional reference materials without large-scale data collection or patient contact. Our dataset follows this approach, generating approximately 560 examples per adapter from a frontier model conditioned on **** Hospital's public patient and family guide. To our knowledge, no prior work has used fine-tuned, age-stratified LoRA adapters for pediatric hospital communication.

\subsection{Virtual Reality in Pediatric Care}

VR-based interventions demonstrably reduce procedural pain and anxiety in pediatric patients, from preoperative preparation to bedside procedures. Existing deployments, however, are largely passive: scripted or pre-rendered distraction content that does not respond to patient queries or adapt to developmental stage, despite hospitalized children frequently having unanswered questions about their environment, staff, and upcoming procedures.

Our work addresses this gap by integrating age-targeted fine-tuned adapters into an interactive VR environment. Unlike passive distraction systems, \textit{VR Hospital Avatar} responds to patient-initiated queries calibrated to the child's developmental stage, combining the anxiety-reduction potential of immersive VR with the informational utility of a conversational agent grounded in hospital-specific knowledge.

\section{Dataset}\label{sec3}

The dataset used to train the models for this project was generated by Claude Opus 4.6 in extended thinking mode. Claude was provided with the Children's Hospital Patient and Family Guide \cite{hospitalguide}, which contains information about the hospital, its routines, procedures, and answers questions that new patients and caretakers may have about their time at the hospital.

From this guide, question and answer pairs were created for six different categories. These categories include Emotional Reassurance, FAQs/General Curiosity, Hospital Rules and Routines, What to Expect, ``Who Are These People?'', and Edge Cases. Questions and age-appropriate answers were generated for each of the different age groups, ages 5-11 and 12-18.

The five primary categories were each targeted at 100 question-and-answer pairs per age group (200 per category). Edge Cases (covering out-of-scope requests, safety-boundary probes, distress escalation, meta/identity questions, and boredom or disengagement) was generated separately at a smaller scale and is training-only, with no corresponding validation examples. After a subsequent quality-curation pass, delivered per-category counts vary from the original targets: \texttt{edge\_cases} (50), \texttt{emotional\_reassurance} (225), \texttt{faqs\_general\_curiosity} (213), \texttt{hospital\_rules\_and\_routines} (212), \texttt{what\_to\_expect} (218), and \texttt{who\_are\_these\_people} (205), for 1123 source training examples in total. After the role-based train/validation split described in Section~\ref{sec:methodology}, this yields 562 training examples for the \texttt{age\_5\_11} adapter and 561 for the \texttt{age\_12\_18} adapter.

These 1123 examples are placed in six separate .jsonl files, one per category.

These examples are formatted as follows:

\begin{verbatim}
{"instruction": "...", "response": "...",
"role": "5-11|12-18", "category": "..."}
\end{verbatim}

The ``instruction'' key holds the user's question, and the ``response'' key holds the appropriate answer to this question. The ``role'' key denotes the age group of the user asking the question; this will always be ``5-11'' or ``12-18''. Lastly, the ``category'' key holds the category of the question asked, which is one of the six categories mentioned prior.

One important thing to note in this dataset is the lack of a system prompt. This is a deliberate choice when designing both the dataset and the project at large: as this guide will be part of the larger Virtual Hospital project, we must ensure it will fit smoothly with the rest of the project. As the team working on the VR and spatial navigation aspects of the Virtual Hospital already relies on their own system prompt for navigation, refraining from including our own system prompt avoids any potential conflicts.

In addition to this, when using a system prompt, we risk the model simply learning the prompt, and becoming overly reliant on it. When the system prompt is removed, this could lead to a degradation in both quality and the needed tone in the model's responses. By avoiding the use of a system prompt in the first place, we can ensure that the model is learning the needed tone in its weights, and not just following the prompt.

\section{Methodology}\label{sec:methodology}

\subsection{System Overview}

This work extends a virtual hospital framework for safe, naturalistic communication between pediatric patients, caregivers, and clinical decision-makers. The platform instantiates a digital recreation of **** Hospital (Location witheld) in which users interact with an AI-powered guide, routed to one of two age-targeted response modules (5--11 or 12--18) determined by the established user profile. Emotional-support queries are in scope only as age-appropriate reassurance and redirection; clinical assessment and crisis intervention are explicitly out of scope. No system prompt is used: the adapter weights encode the age-appropriate register directly, avoiding conflicts with the VR application's own navigation system prompt.

\subsection{Base Models and Infrastructure}

We use the Llama 3.2 Instruct model family \cite{dubey2024llama3} at 1B and 3B scales, loaded in 4-bit quantized form via \texttt{mlx\_lm} \cite{apple2023mlx} on Apple Silicon. All training and evaluation runs on an Apple Mac Mini M4 (16 GB unified memory), a deliberate constraint modeling deployment in resource-limited clinical settings without cloud infrastructure. Full-parameter fine-tuning of the 3B model ($\sim$30 GB in bfloat16 before activations) exceeds this ceiling by nearly 2$\times$; the 1B model is borderline with no headroom. LoRA is therefore the only viable fine-tuning strategy on this hardware for both scales.

\subsection{LoRA as Declarative Programming}

We employ Low-Rank Adaptation (LoRA) \cite{hu2022lora} to specialize the base models without modifying full model weights. A pre-trained weight matrix $W_0 \in \mathbb{R}^{d \times k}$ is held frozen; a low-rank update is expressed as:

\begin{equation}
\Delta W = BA
\end{equation}

where $B \in \mathbb{R}^{d \times r}$ and $A \in \mathbb{R}^{r \times k}$, with rank $r \ll \min(d, k)$. We insert adapters into all linear projection modules within the targeted transformer layers, using \texttt{scale = 20.0} held constant across all runs. In the \texttt{mlx\_lm} framework, \texttt{scale} is a raw output multiplier applied after the low-rank update (not the alpha/rank ratio used in some other implementations), so scale~=~20.0 applies identically to both configurations regardless of rank.

\textbf{Adapter configurations evaluated.} Two configurations were trained and compared:

\begin{itemize}
    \item \textbf{Standard adapter} (rank $r = 8$, \texttt{num\_layers = 16}): covers the upper 16 transformer layers, 57\% of the 3B model and 100\% of the 1B model.
    \item \textbf{Fast adapter} (rank $r = 4$, \texttt{num\_layers = 8}): covers the upper 8 layers, 29\% of the 3B model and 50\% of the 1B model. Prioritizes efficiency: fewer trainable parameters and faster inference.
\end{itemize}

The two configurations differ in both rank and layer count simultaneously, making it impossible to attribute performance differences to either factor in isolation; this is a known limitation of the comparison design.

\subsection{Synthetic Data Generation}

Following the LIMA hypothesis \cite{zhou2023lima}, which holds that a small number of carefully curated examples suffices for alignment, we construct synthetic training corpora rather than collecting naturalistic patient data, motivated both by the impracticality of obtaining real pediatric clinical dialogues at scale and by ethical constraints. Training data was generated by Claude Opus 4.6 (Anthropic; March 2026) with **** Hospital's publicly available patient and family reference materials attached as grounding artifacts. Data is organized into six topic categories:

\begin{itemize}
    \item \textbf{What to Expect}: medical procedures, equipment, and sensory experiences
    \item \textbf{Who Are These People?}: age-appropriate explanations of clinical staff roles
    \item \textbf{Hospital Rules \& Routines}: visiting hours, daily schedules, and preparation
    \item \textbf{Emotional Reassurance}: responses to fear or confusion (non-diagnostic)
    \item \textbf{FAQs / General Curiosity}: open-ended questions a child or caregiver might ask
    \item \textbf{Edge Cases}: out-of-scope requests, safety-boundary probes, distress escalation, meta/identity questions, and boredom or disengagement; training-only, generated in a follow-up pass
\end{itemize}

Each adapter serves a single age group, with approximately 560 training examples per adapter (562 for \texttt{age\_5\_11}, 561 for \texttt{age\_12\_18}) across six categories, totaling 1123 source examples. Hard safety constraints were enforced at generation time: the model was instructed never to imply a diagnosis, never to recommend medications, and to redirect emergency questions to clinical staff. A held-out validation set of 50 examples per age group was generated in a separate pass, drawn from the five primary categories only; \texttt{edge\_cases} has no validation counterpart. Because both sets share the same generative process, perplexity measures in-distribution coherence rather than out-of-distribution generalization.

\subsection{Training Protocol}

Adapters are trained using the Adam optimizer with a cosine decay learning rate schedule: peak learning rate $1 \times 10^{-5}$, linear warmup over the first 100 steps, decaying to zero over 600 total steps. Batch size is 4; maximum sequence length is 768 tokens; LoRA dropout is 0.0. Loss is computed only over assistant response tokens (\texttt{mask\_prompt: true}). With approximately 560 training examples and batch size 4, 600 steps corresponds to approximately 4 passes over the training data. Both configurations are trained across both model sizes and both age groups, yielding 8 total runs (2 configurations $\times$ 2 model sizes $\times$ 2 age groups), all using seed 42.

\subsection{Evaluation Framework}

Three categories of evaluation metrics are reported. \textbf{Readability:} five formula-based metrics are computed over model-generated responses using the \texttt{textstat} library \cite{bansal2023textstat}: Flesch-Kincaid grade level \cite{kincaid1975fk} (primary metric; pass/fail threshold FK~$\leq$~7.0 for 5--11 outputs), SMOG \cite{mclaughlin1969smog} (health-communication-specific, unreliable on very short texts), Gunning Fog \cite{gunning1952fog} (penalises polysyllabic words), Coleman-Liau \cite{coleman1975cli} (character-based cross-check), and type-token ratio (lexical diversity, descriptive only). \textbf{Latency and throughput:} per-response latency and tokens-per-second are measured under sequential, single-request conditions on the Mac Mini M4 (no batching, matching VR deployment conditions); a 1.0-second average threshold reflects real-time interaction requirements. \textbf{Inter-role style separation:} a TF-IDF~$+$ logistic regression classifier (scikit-learn \cite{pedregosa2011sklearn}, max 5{,}000 features, unigram+bigram, max 1{,}000 LR iterations) is trained on the combined 50-prompt validation outputs of both adapters using 5-fold cross-validation; accuracy~$\geq$~0.90 indicates strong register separation (chance~=~0.50).

\section{Results}\label{sec5}

This section reports results for the eight training runs described in Section 4.5: two adapter configurations (Standard, $r = 8$, \texttt{num\_layers = 16}; Fast, $r = 4$, \texttt{num\_layers = 8}) across two model sizes (1B, 3B) and two age groups (5--11, 12--18). All results correspond to the final checkpoint at step 600. Base model (no-adapter) results are included as a zero-shot reference.

\subsection{Training Dynamics and Perplexity}

Table~\ref{tab:res1} reports post-hoc validation perplexity for all eight Llama adapter configurations. Values are computed via a full-sequence cross-entropy pass over the validation set (50 examples per age group) using the step-600 checkpoint; the loss includes all tokens, unlike the masked training objective. All eight adapters converge without divergence.

The perplexity results mirror the configuration-ordering crossover observed in readability. For the 3B base model, Standard achieves lower perplexity than Fast on both age groups (5--11: 18.17 vs.\ 20.77; 12--18: 16.95 vs.\ 21.61), indicating a better-fit adapter. For the 1B base model, the ordering reverses: Fast achieves lower perplexity than Standard (5--11: 22.32 vs.\ 24.17; 12--18: 19.79 vs.\ 22.31). This corroborates the FK $\leq 7.0$ pass-rate crossover and strengthens the capacity-regularization interpretation: Standard (higher rank) overfits relative to the 1B model's capacity, while Fast (lower rank) underfits relative to the 3B model.

\begin{table}[!t]
  \centering
  \caption{Validation perplexity at step 600 for all Llama adapter configurations (post-hoc full-sequence loss pass).}
  \label{tab:res1}
  \small
  \begin{tabular}{llcc}
    \hline
    Config. & Age & Val Loss & PPL \\
    \hline
    Standard-3B & 5--11  & 2.900 & 18.17 \\
    Standard-3B & 12--18 & 2.830 & 16.95 \\
    Standard-1B & 5--11  & 3.185 & 24.17 \\
    Standard-1B & 12--18 & 3.105 & 22.31 \\
    Fast-3B     & 5--11  & 3.034 & 20.77 \\
    Fast-3B     & 12--18 & 3.073 & 21.61 \\
    Fast-1B     & 5--11  & 3.106 & 22.32 \\
    Fast-1B     & 12--18 & 2.985 & 19.79 \\
    \hline
  \end{tabular}
\end{table}

\subsection{Readability}

Table~\ref{tab:res2} reports readability metrics for all six configurations. All four fine-tuned adapters substantially outperform the base models on the FK $\leq 7.0$ target for the 5--11 group (72--84\% pass rate vs.\ 12--14\% for base models), driven largely by shorter response length (72--99 avg words vs.\ 152--210 for base models). The 12--18 adapters produce consistently higher FK grade levels than the 5--11 adapters (gap: 2.4--3.2 FK points), confirming that the two adapters have learned distinct communication registers.

A \textbf{configuration-ordering crossover} is visible in Table~\ref{tab:res2}: for the Standard configuration, Standard-3B outperforms Standard-1B on FK $\leq 7.0$ pass rate (84\% vs.\ 72\%); for the Fast configuration, this ordering reverses: Fast-1B outperforms Fast-3B (82\% vs.\ 76\%). This crossover is consistent across all five readability metrics for the 5--11 group and is the central empirical observation of this evaluation. For the four grade-level metrics (FK, SMOG, Gunning Fog, Coleman-Liau), lower values indicate greater age-appropriateness for the 5--11 group; for TTR, higher values indicate more lexically varied responses, which is preferable. The crossover is directionally consistent under both interpretations.

\begin{table*}[!t]
  \caption{Readability metrics (outputs from age-matched held-out prompts, 50 examples per group).}
  \label{tab:res2}
  \small
  \begin{tabular}{llccccccc}
    \hline
    Configuration & Age Group & FK Grade & SMOG & Gunning Fog & Coleman-Liau & TTR & Avg Words & FK $\leq$ 7.0 \\
    \hline
    Standard-1B & 5--11  & 6.1  & 8.8  & 7.7  & 6.7  & 0.712 & 74  & 36/50 (72\%) \\
    Standard-1B & 12--18 & 8.8  & 11.7 & 11.3 & 10.3 & 0.701 & 90  & ---          \\
    Standard-3B & 5--11  & 5.5  & 8.5  & 7.4  & 6.3  & 0.724 & 75  & 42/50 (84\%) \\
    Standard-3B & 12--18 & 8.7  & 11.4 & 10.9 & 10.0 & 0.714 & 92  & ---          \\
    Fast-1B     & 5--11  & 5.5  & 8.4  & 7.2  & 6.2  & 0.701 & 72  & 41/50 (82\%) \\
    Fast-1B     & 12--18 & 8.4  & 11.3 & 10.6 & 9.4  & 0.650 & 92  & ---          \\
    Fast-3B     & 5--11  & 5.9  & 8.9  & 7.9  & 6.6  & 0.636 & 99  & 38/50 (76\%) \\
    Fast-3B     & 12--18 & 8.3  & 11.2 & 10.7 & 9.7  & 0.613 & 122 & ---          \\
    Base-1B     & 5--11  & 9.5  & 12.1 & 12.2 & 9.0  & 0.508 & 201 & 7/50 (14\%)  \\
    Base-1B     & 12--18 & 10.8 & 13.1 & 13.4 & 11.2 & 0.527 & 210 & ---          \\
    Base-3B     & 5--11  & 9.4  & 12.1 & 12.2 & 9.6  & 0.609 & 152 & 6/50 (12\%)  \\
    Base-3B     & 12--18 & 10.6 & 12.9 & 13.3 & 11.4 & 0.565 & 206 & ---          \\
    \hline
  \end{tabular}
\end{table*}

\subsection{Latency and Throughput}

Table~\ref{tab:res3} reports inference latency for all configurations (sequential, no batching, Mac Mini M4). \textbf{Fast-1B is the only configuration meeting the 1.0-second average latency target} for the 5--11 adapter (0.93s avg, 70\% of responses under target). The Fast-1B 12--18 adapter averages 1.20s, missing the threshold. All 3B configurations exceed 2.3 seconds on average. Fast-1B's latency advantage over Base-1B is driven primarily by output length rather than throughput: Fast-1B generates 88 avg tokens at 94.3 tok/s, while Base-1B generates 247 avg tokens at 122.3 tok/s; the base model is faster per token but produces responses nearly 3$\times$ longer, making fine-tuning the dominant factor in meeting the latency target. Counterintuitively, Fast-3B is slower than Standard-3B despite fewer adapter parameters: Fast-3B generates $\sim$30\% more tokens per response (120/153 vs.\ 92/116 avg tokens), overriding any throughput benefit from the smaller adapter.

\begin{table*}[!t]
  \caption{Inference latency and throughput (50 prompts per age group, sequential, Mac Mini M4).}
  \label{tab:res3}
  \small
  \begin{tabular}{llcccccc}
    \hline
    Config. & Age & Avg Lat. (s) & Min (s) & Max (s) & Avg Tok & Tok/s & $<$1.0s \\
    \hline
    Standard-1B & 5--11  & 1.09 & 0.66 & 1.69 & 89  & 81.7  & 21/50 (42\%) \\
    Standard-1B & 12--18 & 1.36 & 0.62 & 2.02 & 114 & 83.7  & 3/50 (6\%)   \\
    Standard-3B & 5--11  & 2.37 & 1.31 & 3.45 & 92  & 38.8  & 0/50 (0\%)   \\
    Standard-3B & 12--18 & 2.94 & 1.75 & 5.43 & 116 & 39.3  & 0/50 (0\%)   \\
    Fast-1B     & 5--11  & 0.93 & 0.50 & 1.40 & 88  & 94.3  & 35/50 (70\%) \\
    Fast-1B     & 12--18 & 1.20 & 0.75 & 1.79 & 115 & 95.4  & 12/50 (24\%) \\
    Fast-3B     & 5--11  & 2.83 & 1.31 & 6.80 & 120 & 41.9  & 0/50 (0\%)   \\
    Fast-3B     & 12--18 & 3.63 & 0.99 & 7.33 & 153 & 41.7  & 1/50 (2\%)   \\
    Base-1B     & 5--11  & 2.01 & 0.57 & 2.48 & 247 & 122.3 & 2/50 (4\%)   \\
    Base-1B     & 12--18 & 2.14 & 0.25 & 2.46 & 265 & 121.8 & 4/50 (8\%)   \\
    Base-3B     & 5--11  & 3.99 & 0.66 & 6.38 & 190 & 46.5  & 3/50 (6\%)   \\
    Base-3B     & 12--18 & 5.39 & 0.67 & 6.29 & 264 & 48.7  & 1/50 (2\%)   \\
    \hline
  \end{tabular}
\end{table*}

\subsection{Inter-Role Style Separation}

Table~\ref{tab:res4} reports TF-IDF + logistic regression classification accuracy (5-fold CV; chance = 0.50) for distinguishing 5--11 vs.\ 12--18 adapter outputs. All fine-tuned configurations achieve strong separation (0.89--0.94); base models reach only moderate separation (0.66--0.70), confirming that fine-tuning is the primary driver of distinct age registers. The same crossover observed in readability is present here: Standard-3B outperforms Standard-1B (0.920 vs.\ 0.890), while Fast-1B outperforms Fast-3B (0.940 vs.\ 0.900).

\begin{table}[!t]
  \centering
  \caption{Inter-role TF-IDF + LR classification accuracy (5-fold CV; chance = 0.50).}
  \label{tab:res4}
  \small
  \begin{tabular}{lc}
    \hline
    Configuration & Classification Accuracy \\
    \hline
    Standard-1B & $0.890 \pm 0.102$ \\
    Standard-3B & $0.920 \pm 0.051$ \\
    Fast-1B     & $0.940 \pm 0.073$ \\
    Fast-3B     & $0.900 \pm 0.077$ \\
    Base-1B     & $0.660 \pm 0.097$ \\
    Base-3B     & $0.700 \pm 0.063$ \\
    \hline
  \end{tabular}
\end{table}


\subsection{Discussion}

\textbf{Fast-1B is the recommended configuration} for VR deployment: it is the only configuration meeting the 1.0-second latency target (0.93s avg) while maintaining strong readability (FK 5.5 avg) and style separation (0.940 accuracy). The similarity between Standard and Fast on most quality metrics suggests that register adaptation for this task is achievable at rank 4 with 8 adapted layers: high-quality, constrained training data may reduce the adapter capacity required to produce measurable style change.

The \textbf{crossover observation} (Standard adapter better on 3B, Fast adapter better on 1B for both readability and style separation) is the central empirical finding. We tentatively interpret this as a capacity-regularization effect: the 1B model's limited capacity may benefit from the implicit regularization of a smaller adapter, while the 3B model can exploit the extra degrees of freedom of rank 8. Since rank and layer count co-vary here and all runs use a single seed, the crossover is not yet statistically confirmed, motivating the controlled rank sweep in Section~\ref{sec:ranksweep}.

Structural limitations are discussed fully in Section~\ref{sec:limitations}.

\section{Extended Experiments}\label{sec:extended}

Following the evaluation of the original eight runs, two follow-up experiments were conducted to address the structural limitations identified in Section~\ref{sec:limitations}: a controlled rank sweep to isolate the crossover mechanism, and cross-architecture evaluation on two additional model families.

\subsection{Rank Sweep: Mechanism Identification}\label{sec:ranksweep}

The Standard and Fast configurations differ in both \texttt{rank} and \texttt{num\_layers} simultaneously, making it impossible from the original eight runs to attribute the crossover to either factor. To isolate rank, we fixed \texttt{num\_layers~=~8} and swept $r \in \{2, 4, 8, 16\}$ across four architectures (Llama~1B, Llama~3B, Qwen~2 0.5B, SmolLM2~360M) with two seeds (42 and 1337), targeting the \texttt{age\_5\_11} role only. This yields 32 training runs.

Table~\ref{tab:ranksweep} reports FK~$\leq$~7.0 pass rate as a function of rank averaged across seeds. The four models exhibit two distinct profiles:

\begin{itemize}
    \item \textbf{Decreasing (small-model regime):} Llama~1B peaks at $r=2$ (81\%) and degrades monotonically at higher rank. Qwen~0.5B likewise peaks at $r=2$ (78\%).
    \item \textbf{Increasing (large-model regime):} Llama~3B improves through $r=8$ (69\%) before plateauing. SmolLM2~360M improves monotonically through $r=16$ (76\%) and incurs the largest penalty at $r=2$ (49\%).
\end{itemize}

Because \texttt{num\_layers} is held constant across all sweep runs, the crossover between $r=4$ and $r=8$ cannot be attributed to layer coverage. \textbf{Rank is the operative variable}, consistent with a capacity-regularization mechanism: models with limited representational capacity benefit from the implicit constraint of lower rank; deeper models can exploit additional rank for stronger style adaptation.

Notably, SmolLM2~360M falls in the large-model regime despite having fewer total parameters than any Llama configuration. Its 32 transformer layers and hidden dimension of 960 provide per-layer capacity comparable to deeper models, suggesting that \textbf{transformer depth, not total parameter count, determines rank sensitivity}.

\begin{table}[!t]
  \centering
  \caption{Rank sweep FK $\leq 7.0$ pass rate (\%, mean $\pm$ std across two seeds). \texttt{num\_layers~=~8} fixed. Best per model in \textbf{bold}.}
  \label{tab:ranksweep}
  \small
  \begin{tabular}{lcccc}
    \hline
    Model        & $r=2$               & $r=4$     & $r=8$               & $r=16$ \\
    \hline
    Llama 1B     & \textbf{81 $\pm$ 1} & 65 $\pm$ 1 & 63 $\pm$ 5         & 63 $\pm$ 3 \\
    Llama 3B     & 56 $\pm$ 4          & 60 $\pm$ 4 & \textbf{69 $\pm$ 9} & 68 $\pm$ 0 \\
    Qwen 0.5B    & \textbf{78 $\pm$ 0} & 72 $\pm$ 4 & 75 $\pm$ 3         & 68 $\pm$ 2 \\
    SmolLM2 360M & 49 $\pm$ 7          & 69 $\pm$ 5 & 70 $\pm$ 6         & \textbf{76 $\pm$ 8} \\
    \hline
  \end{tabular}
\end{table}

\subsection{Cross-Architecture Validation}\label{sec:crossarch}

To test whether the crossover generalizes beyond Llama, we evaluate Qwen~2 0.5B (4-bit quantized, ChatML template) and SmolLM2 360M \cite{allal2025smollm2} (4-bit quantized, ChatML template) under the Standard ($r=8$, \texttt{num\_layers~=~16}) and Fast ($r=4$, \texttt{num\_layers~=~8}) configurations, both age groups, seed 42. SmolLM2's chat template injects a default system message when no system role is present; training data was generated with empty-content system messages to suppress this, and inference uses a matching empty system context to maintain consistency with training conditions.

Tables~\ref{tab:crossarch} and~\ref{tab:crossarchclf} summarize readability, latency, and classifier results. Several findings follow:

\textbf{Qwen 0.5B: Standard wins on FK, Fast wins on classifier.} This split is resolved by the rank sweep (Table~\ref{tab:ranksweep}): Qwen's monotonically decreasing rank curve places it in the small-model regime. Both configurations clear the 1.0s latency target (100\% and 98\% of age~5--11 responses, respectively) due to high token throughput ($\sim$175 tok/s).

\textbf{SmolLM2 Standard dominates on all metrics.} It achieves 84\% FK pass rate and 0.950 classifier accuracy, matching Standard-3B on FK and ranking second overall on the classifier, while averaging 0.81\,s for the age~5--11 role (88\% under target). The Fast adapter, by contrast, produces substantially longer responses (111 vs.\ 76 avg tokens) and falls to 64\% FK pass rate, indicating rank~4 is under-parameterized for its 32-layer architecture. SmolLM2 Standard represents a strong quality-latency trade-off: the depth of a large model with the throughput of a small one.

This raises four viable deployment candidates spanning the latency-quality trade-off; the full ordering, updated after the Qwen 2.5 sweep, is given at the end of Section~\ref{sec:qwen25}.

\begin{table}[!t]
  \centering
  \caption{Cross-architecture readability and latency for the age~5--11 role (50 examples, seed 42).}
  \label{tab:crossarch}
  \small
  \begin{tabular}{llccccc}
    \hline
    Config. & FK & FK $\leq$7.0 & Avg Lat. & $<$1.0s & Avg Tok \\
    \hline
    Qwen Std.    & 6.2 & 37/50 (74\%) & 0.59\,s & 50/50 (100\%) & 88 \\
    Qwen Fast    & 5.8 & 34/50 (68\%) & 0.46\,s & 49/50 (98\%)  & 83 \\
    SmolLM2 Std. & 5.8 & 42/50 (84\%) & 0.81\,s & 44/50 (88\%)  & 76 \\
    SmolLM2 Fast & 6.1 & 32/50 (64\%) & 1.08\,s & 34/50 (68\%)  & 111 \\
    \hline
  \end{tabular}
\end{table}

\begin{table}[!t]
  \centering
  \caption{Inter-role classifier accuracy for cross-architecture configurations (5-fold CV; chance~=~0.50).}
  \label{tab:crossarchclf}
  \small
  \begin{tabular}{lc}
    \hline
    Configuration    & Accuracy \\
    \hline
    Qwen Standard    & $0.940 \pm 0.049$ \\
    Qwen Fast        & $0.960 \pm 0.058$ \\
    SmolLM2 Standard & $0.950 \pm 0.032$ \\
    SmolLM2 Fast     & $0.920 \pm 0.040$ \\
    \hline
  \end{tabular}
\end{table}

\subsection{Qwen 2.5 Family Sweep}\label{sec:qwen25}

To test whether the rank regime prediction extends across a capacity ladder within a single family, we trained and evaluated the Qwen 2.5 Instruct family at 0.5B, 1.5B, and 3B under both Standard ($r=8$, \texttt{num\_layers~=~16}) and Fast ($r=4$, \texttt{num\_layers~=~8}) configurations, both age groups, seed 42. This yields 12 training runs. Results are reported in Tables~\ref{tab:qwen25} and~\ref{tab:qwen25clf}.

\textbf{Fast outperforms Standard at every size; no crossover.} FK $\leq$ 7.0 pass rates are 76\% vs.\ 72\% (0.5B), 70\% vs.\ 48\% (1.5B), and 76\% vs.\ 58\% (3B) for Fast vs.\ Standard, respectively. This contrasts with the Llama results, where Standard wins on the 3B model. The uniform Fast advantage is consistent with the rank sweep placing Qwen in the small-model regime (peak FK at $r=2$): Standard ($r=8$) is over-parameterized across the entire Qwen 2.5 capacity ladder. The Standard penalty worsens with size: the 1.5B Standard achieves only 48\% FK pass rate, the weakest result of any fine-tuned adapter in this evaluation.

\textbf{Throughput degrades steeply with scale.} Tokens per second fall from $\sim$178 (0.5B Fast) to $\sim$75 (1.5B Fast) to $\sim$44 (3B Fast). This decline is much steeper than the Llama 1B-to-3B degradation, likely attributable to architecture differences in the Qwen 2.5 attention implementation under 4-bit quantization on Apple Silicon. As a consequence, only the 0.5B Fast meets the 1.0-second latency target (50/50 under target); the 1.5B Fast averages 0.98\,s with 26/50 responses under target.

\textbf{Qwen 2.5 1.5B Fast sets the highest classifier accuracy of any adapter evaluated (0.980 $\pm$ 0.024).} This exceeds the previous best (Qwen 2 Fast, 0.960) by two points. The result suggests that a 1.5B model in the small-model regime, where rank 4 is the better configuration, can encode stronger register contrast than any Llama or Qwen 2 configuration, at the cost of borderline latency.

\textbf{Updated deployment recommendation.} Incorporating the Qwen 2.5 results, the full priority ordering is:
\begin{itemize}
    \item \textit{Lowest latency:} Qwen 2 Fast or Qwen 2.5-0.5B Fast (0.46\,s avg; 100\% under 1.0\,s; 0.960 and 0.950 classifier, respectively).
    \item \textit{Best quality-latency balance:} SmolLM2 Standard (84\% FK; 0.950 classifier; 0.81\,s avg).
    \item \textit{Highest style separation:} Qwen 2.5-1.5B Fast (0.980 classifier; 70\% FK; 0.98\,s avg; borderline on latency target).
    \item \textit{Llama-only deployment:} Fast-1B (0.93\,s avg; 82\% FK; 0.940 classifier).
\end{itemize}

\begin{table}[!t]
  \centering
  \caption{Qwen 2.5 family readability and latency for the age~5--11 role (50 examples, seed 42).}
  \label{tab:qwen25}
  \small
  \begin{tabular}{lccccc}
    \hline
    Config. & FK & FK $\leq$7.0 & Avg Lat. & $<$1.0s & Tok/s \\
    \hline
    Qwen2.5-0.5B Fast & 5.8 & 38/50 (76\%) & 0.46\,s & 50/50 (100\%) & 177.7 \\
    Qwen2.5-0.5B Std. & 6.1 & 36/50 (72\%) & 0.57\,s & 50/50 (100\%) & 151.8 \\
    Qwen2.5-1.5B Fast & 6.2 & 35/50 (70\%) & 0.98\,s & 26/50 (52\%)  & 75.2  \\
    Qwen2.5-1.5B Std. & 7.0 & 24/50 (48\%) & 1.31\,s & 6/50 (12\%)   & 68.2  \\
    Qwen2.5-3B Fast   & 5.9 & 38/50 (76\%) & 1.89\,s & 2/50 (4\%)    & 43.8  \\
    Qwen2.5-3B Std.   & 6.8 & 29/50 (58\%) & 2.23\,s & 0/50 (0\%)    & 40.1  \\
    \hline
  \end{tabular}
\end{table}

\begin{table}[!t]
  \centering
  \caption{Inter-role classifier accuracy for Qwen 2.5 configurations (5-fold CV; chance~=~0.50).}
  \label{tab:qwen25clf}
  \small
  \begin{tabular}{lc}
    \hline
    Configuration     & Accuracy \\
    \hline
    Qwen2.5-0.5B Fast & $0.950 \pm 0.032$ \\
    Qwen2.5-0.5B Std. & $0.940 \pm 0.097$ \\
    Qwen2.5-1.5B Fast & $\mathbf{0.980 \pm 0.024}$ \\
    Qwen2.5-1.5B Std. & $0.890 \pm 0.073$ \\
    Qwen2.5-3B Fast   & $0.970 \pm 0.040$ \\
    Qwen2.5-3B Std.   & $0.930 \pm 0.068$ \\
    \hline
  \end{tabular}
\end{table}

\subsection{Layer Sweep: Independent Effect of \texttt{num\_layers}}\label{sec:layersweep}

The rank sweep held \texttt{num\_layers} fixed at 8 to isolate rank, confirming rank as the operative crossover factor. However, the original Standard and Fast configurations also differ in depth (16 vs.\ 8 layers), leaving open whether layer count contributes independently. To quantify this, we fixed \texttt{rank=4} and swept \texttt{num\_layers} $\in \{4, 8, 16\}$ across Llama 1B and 3B with two seeds (42 and 1337), targeting \texttt{age\_5\_11}. This yields 12 runs; the \texttt{num\_layers=8} configurations reuse existing rank sweep adapters. Results are in Table~\ref{tab:layersweep}.

\textbf{\texttt{num\_layers} is an independent secondary contributor.} At fixed rank=4, layers=16 outperforms layers=8 by $+4.0\%$ (1B: 71\% vs.\ 67\%) and $+7.0\%$ (3B: 65\% vs.\ 58\%). The Standard configuration's depth advantage ($16$ vs.\ $8$ layers) thus contributes to its superiority on 3B beyond rank alone. The original crossover reflects the joint effect of both rank and depth, with rank identified as the dominant factor by the rank sweep.

\textbf{Small models prefer fewer layers.} On 1B, the peak FK pass rate occurs at \texttt{num\_layers=4} (75.0\%), declining to 67.0\% at layers=8 and recovering partially to 71.0\% at layers=16. This is consistent with capacity over-parameterization: adapting more layers with an already-constrained model spreads the adapter's capacity too thin. On 3B, layers=16 is the clear winner (65.0\%), consistent with larger models benefiting from broader coverage.

\textbf{Latency scales with layer count.} For 1B, average latency rises from 0.92\,s (layers=4) to 0.93\,s (layers=8) to 1.07\,s (layers=16). The layers=16 configuration marginally exceeds the 1.0-second real-time target; layers=4 achieves the best FK and lowest latency simultaneously, making it the recommended configuration for latency-critical 1B deployments.

\begin{table}[!t]
  \centering
  \caption{Layer sweep at fixed rank=4, role \texttt{age\_5\_11} (mean $\pm$ std across seeds 42 and 1337).}
  \label{tab:layersweep}
  \small
  \begin{tabular}{lcccc}
    \hline
    Config. & FK avg & FK $\leq$7.0 & Avg Lat. & Tok/s \\
    \hline
    1B, layers=4  & $6.16 \pm 0.10$ & $75.0\% \pm 5.0$ & $0.92\,\text{s} \pm 0.00$ & $101.6 \pm 1.0$ \\
    1B, layers=8  & $6.28 \pm 0.02$ & $67.0\% \pm 1.0$ & $0.93\,\text{s} \pm 0.00$ & $93.5 \pm 1.0$  \\
    1B, layers=16 & $6.46 \pm 0.13$ & $71.0\% \pm 3.0$ & $1.07\,\text{s} \pm 0.03$ & $83.4 \pm 0.4$  \\
    3B, layers=4  & $6.46 \pm 0.12$ & $58.0\% \pm 6.0$ & $2.45\,\text{s} \pm 0.11$ & $44.5 \pm 0.1$  \\
    3B, layers=8  & $6.62 \pm 0.11$ & $58.0\% \pm 2.0$ & $2.59\,\text{s} \pm 0.03$ & $42.2 \pm 0.0$  \\
    3B, layers=16 & $6.52 \pm 0.01$ & $65.0\% \pm 1.0$ & $2.20\,\text{s} \pm 0.05$ & $38.4 \pm 0.3$  \\
    \hline
  \end{tabular}
\end{table}

\section{Limitations}\label{sec:limitations}

\textbf{Model family and scale.} The original eight runs use Llama 3.2 at 1B and 3B only. Extended experiments (Section~\ref{sec:extended}) add Qwen~2 0.5B, SmolLM2 360M, and Qwen 2.5 (0.5B, 1.5B, 3B), confirming the crossover is not Llama-specific. Models above 3B and the Gemma and Mistral families remain untested.

\textbf{Confounded adapter configurations (resolved).} Standard and Fast differ in both rank and \texttt{num\_layers}. The rank sweep (Section~\ref{sec:ranksweep}) confirms rank as the operative factor; the layer sweep (Section~\ref{sec:layersweep}) confirms depth as an independent secondary contributor. The crossover reflects the joint effect of both, with rank dominant.

\textbf{Single-seed runs (partially resolved).} The original eight runs use seed 42 only. The rank and layer sweeps use two seeds and report mean $\pm$ std; variance for the cross-architecture runs is not available.

\textbf{Synthetic data and in-distribution evaluation.} Training and validation data share the same generative process and topic distribution, so perplexity measures in-distribution coherence only. Human-written references and external evaluators are needed to confirm genuinely age-appropriate output rather than learned stylistic imitation.

\textbf{No clinical validation.} The system has not been tested with pediatric patients, caregivers, or clinical staff. \textit{VR Hospital Avatar} is a research prototype requiring clinical usability studies and institutional review before deployment.

\textbf{Safety evaluation.} No automated safety evaluation was performed. A guard model for output filtering is planned for future deployment; the current system relies on the Llama 3.2 Instruct base model's safety alignment and the safety constraints enforced at data generation time.

\section{Conclusion}

This paper demonstrates that LIMA-style synthetic fine-tuning on consumer hardware can produce measurable, age-appropriate communication behavior in a pediatric hospital context. Fine-tuned adapters substantially outperform zero-shot base models across all measured dimensions: readability pass rates improve from 12--14\% to 72--84\% on the FK $\leq 7.0$ target for the 5--11 group; inter-role style separation improves from 0.66--0.70 to 0.89--0.94 classification accuracy; and the Fast-1B configuration reduces average latency from 2.01s (Base-1B) to 0.93s.

The central finding is a \textbf{configuration-ordering crossover}: Standard ($r=8$, 16 layers) outperforms Fast on the 3B model, while Fast ($r=4$, 8 layers) outperforms on 1B, consistent across all five readability metrics and the classification task. The rank sweep (Section~\ref{sec:ranksweep}) confirms \textbf{rank as the dominant factor} via capacity regularization, with transformer depth, not parameter count, predicting the optimal rank regime; the layer sweep (Section~\ref{sec:layersweep}) confirms depth as an independent secondary contributor (layers=16 beats layers=8 by 4--7\% FK at fixed rank). Cross-architecture evaluation (Qwen~2 0.5B, SmolLM2 360M, Qwen 2.5 0.5B--3B) confirms the crossover is not Llama-specific; within Qwen 2.5, Fast dominates Standard uniformly, consistent with the small-model rank regime, with Qwen 2.5 1.5B Fast achieving the highest classifier accuracy of any adapter evaluated (0.980).

For deployment, Section~\ref{sec:qwen25} gives the full priority ordering across latency, quality, and style-separation trade-offs; Fast-1B remains the recommended choice for Llama-only deployments. Remaining limitations include in-distribution-only evaluation and the absence of clinical validation with pediatric patients or caregivers (Section~\ref{sec:limitations}).


%\backmatter
\bmsubsection*{Author Contributions}

[Author contributions withheld for double-blind review.]

\bmsubsection*{Acknowledgments}
[Acknowledgments withheld for double-blind review.]

\bmsubsection*{Financial Disclosure}

None reported.

\bmsubsection*{Conflicts of Interest}

The authors declare no conflicts of interest.

%\bibliography{wileyNJD-Chicago}

\begin{thebibliography}{99}

\bibitem{zhou2023lima}
C.~Zhou, P.~Liu, P.~Xu, S.~Iyer, J.~Sun, R.~Maddela, X.~Peng, S.~Shrivastava, M.~Lewis, L.~Zettlemoyer, and O.~Levy, ``LIMA: Less Is More for Alignment,'' in \textit{Advances in Neural Information Processing Systems (NeurIPS)}, 2023.

\bibitem{dubey2024llama3}
A.~Dubey et al., ``The Llama 3 Herd of Models,'' \textit{arXiv preprint arXiv:2407.21783}, Meta AI, 2024.

\bibitem{apple2023mlx}
Apple, ``MLX: Efficient and Flexible Machine Learning on Apple Silicon,'' 2023. [Online]. Available: https://github.com/ml-explore/mlx

\bibitem{hu2022lora}
E.~J.~Hu, Y.~Shen, P.~Wallis, Z.~Allen-Zhu, Y.~Li, S.~Wang, L.~Wang, and W.~Chen, ``LoRA: Low-Rank Adaptation of Large Language Models,'' in \textit{Proc. Int. Conf. Learning Representations (ICLR)}, 2022.

\bibitem{kincaid1975fk}
J.~P.~Kincaid, R.~P.~Fishburne, R.~L.~Rogers, and B.~S.~Chissom, ``Derivation of New Readability Formulas for Navy Enlisted Personnel,'' Research Branch Report 8-75, Naval Technical Training Command, 1975.

\bibitem{mclaughlin1969smog}
G.~H.~McLaughlin, ``SMOG Grading --- A New Readability Formula,'' \textit{J. Reading}, vol.~12, no.~8, pp.~639--646, 1969.

\bibitem{gunning1952fog}
R.~Gunning, \textit{The Technique of Clear Writing}. New York, NY, USA: McGraw-Hill, 1952.

\bibitem{coleman1975cli}
M.~Coleman and T.~L.~Liau, ``A Computer Readability Formula Designed for Machine Scoring,'' \textit{J. Applied Psychology}, vol.~60, no.~2, pp.~283--284, 1975.

\bibitem{bansal2023textstat}
S.~Bansal, ``textstat: Python Library to Calculate Readability Statistics of a Text Object'' (v0.7.13), 2023. [Online]. Available: \url{https://pypi.org/project/textstat/}

\bibitem{dettmers2023qlora}
T.~Dettmers, A.~Pagnoni, A.~Holtzman, and L.~Zettlemoyer, ``QLoRA: Efficient Finetuning of Quantized LLMs,'' in \textit{Advances in Neural Information Processing Systems (NeurIPS)}, 2023.

\bibitem{biderman2024lora}
D.~Biderman, J.~G.~Ortiz, J.~Portes, M.~Paul, P.~Greengard, C.~Jennings, D.~King, S.~Havens, V.~Chiley, J.~Frankle, C.~Blakeney, and J.~P.~Cunningham, ``LoRA Learns Less and Forgets Less,'' \textit{Transactions on Machine Learning Research (TMLR)}, 2024.

\bibitem{pedregosa2011sklearn}
F.~Pedregosa et al., ``Scikit-learn: Machine Learning in Python,'' \textit{J. Machine Learning Research}, vol.~12, pp.~2825--2830, 2011.

\bibitem{allal2025smollm2}
L.~B. Allal et al., ``SmolLM2: When Smol Goes Big,'' \textit{arXiv preprint arXiv:2502.05737}, Hugging Face, 2025.

\bibitem{hospitalguide}
“Children's Hospital Patient and Family Guide,” Children's Hospital **** Health Sciences Centre. https://www.****/childrens-hospital/childrens-hospital-patient-and-family-guide (accessed Mar. 26, 2026).

\end{thebibliography}

\bmsubsection*{Supporting Information}

Additional supporting information can be found online in the Supporting Information
section.

\appendix

%% Repository link removed for double-blind review (deanonymizes via GitHub username).
%% TODO: review with colleague whether AECAI-PRiSM 2026 permits a repo link during
%% review before restoring \url{https://github.com/abarros6/small-bear} here.

\subsection{Data Generation Provenance}

All synthetic data was generated using Claude Opus 4.6 with extended thinking mode via Claude.ai Project Spaces (March 2026), with **** Hospital's publicly available patient and family guide attached as a grounding artifact in dedicated, per-category chats to maintain context isolation; the \texttt{edge\_cases} category was added in a later follow-up pass. Category definitions, per-category counts, and the hard safety rules enforced at generation time are given in Section~\ref{sec:methodology}. Training scripts and evaluation code were developed using Claude Code (claude-sonnet-4-6, v2.1.74).

\nocite{*}% Show all bib entries - both cited and uncited; comment this line to view only cited bib entries;




\end{document}
